\subsubsection{Reducción de datos y error de contacto de la superficie parcial}
\label{sec:pie-rey-reduccion}

Sea $T_k$ la temperatura registrada en el punto de calibración $k$ y
$T_0 = 20\ ^\circ\text{C}$ la temperatura de referencia. El valor de referencia de
cada punto es la longitud de la pila de bloques referida a $T_0$; por la
ecuación~\eqref{eq:e-termal}, e incorporando la corrección de certificado
$c_{cert,k}$,

\begin{align}
  VPC_k &= VP_k + c_{cert,k} + \alpha_{BP}\, L_k\,(T_k - T_0),
  \label{eq:vpc}
\end{align}

donde $VP_k$ es el valor nominal de la pila, $c_{cert,k}$ la corrección acumulada de
los bloques que la componen según su certificado,
$\alpha_{BP} = 11.5\times 10^{-6}\ ^\circ\text{C}^{-1}$ el coeficiente de expansión
térmica del acero de los bloques y $L_k \approx VP_k$. De forma análoga, la
indicación del pie de rey se refiere a $T_0$ según

\begin{equation}
  PC_{kj} = X_{kj}\left[1 - \alpha_i (T_k - T_0)\right],
  \label{eq:pc}
\end{equation}

donde $X_{kj}$ es la lectura en la posición $j$ del punto $k$.

El error de indicación de cada posición es la diferencia entre ambas reducciones.
Sustituyendo \eqref{eq:vpc} y \eqref{eq:pc} y desarrollando,

\begin{align}
  D_{kj} &= PC_{kj} - VPC_k \nonumber\\
  &= X_{kj} - X_{kj}\alpha_i (T_k - T_0) - VP_k - c_{cert,k}
     - \alpha_{BP}\,L_k (T_k - T_0). \label{eq:D-desarrollo}
\end{align}

Como $X_{kj} \approx L_k \approx VP_k$ en los términos correctivos y el pie de rey y
los bloques son de acero ($\alpha_i = \alpha_{BP} = \alpha$),

\begin{align*}
  D_{kj} &= X_{kj} - VP_k - c_{cert,k}
     - \alpha\,VP_k (T_k - T_0) + \alpha\,VP_k (T_k - T_0) \\
  &= X_{kj} - VP_k - c_{cert,k},
\end{align*}

es decir

\begin{equation}
  \therefore\quad D_{kj} \approx X_{kj} - VP_k - c_{cert,k},
  \label{eq:D}
\end{equation}

expresada en micrómetros. La corrección térmica se cancela por ser común a ambos
cuerpos, y la corrección de certificado es constante dentro de cada punto, de modo
que tampoco interviene en la dispersión: el error de contacto que se calcula a
continuación es independiente de $c_{cert,k}$ y de la corrección térmica.

La Tabla~\ref{tab:pie-rey-vpc} presenta, para cada punto, el valor patrón corregido,
el promedio corregido de las tres lecturas y el error de indicación resultante. Al no
disponer de las desviaciones certificadas de los bloques, se adopta $c_{cert,k} = 0$;
esta hipótesis no afecta al cálculo de $E$, pero sí al error de indicación reportado
en cada punto.

\begin{table}[H]
\centering
\caption{Valor patrón corregido, promedio corregido y error de indicación en cada punto de calibración.}
\label{tab:pie-rey-vpc}
\begin{tabular}{rrrr}
\toprule
$VP$ (mm) & $VPC$ (mm) & $PC$ (mm) & Error de indicación ($\mu$m) \\
\midrule
50.00 & 50.001 & 49.999 & -1.7 \\
100.00 & 100.002 & 100.015 & 13.2 \\
150.00 & 150.003 & 150.031 & 28.2 \\
200.00 & 200.003 & 200.030 & 26.4 \\
250.00 & 250.004 & 249.996 & -8.6 \\
\bottomrule
\end{tabular}
\end{table}


\subsubsection*{Error de contacto de la superficie parcial en medición de exteriores}

Sea $E_k$ la dispersión de las tres posiciones en el punto $k$. El error de contacto
se define como la mayor de las dispersiones observadas,

\begin{equation}
  E = \max_k \left( \max_j D_{kj} - \min_j D_{kj} \right),
  \label{eq:E}
\end{equation}

cuyos valores se recogen en la Tabla~\ref{tab:pie-rey-diferencias}.

\begin{table}[H]
\centering
\caption{Diferencias de indicación por posición y dispersión $E_k$ en cada punto de calibración.}
\label{tab:pie-rey-diferencias}
\begin{tabular}{rrrrr}
\toprule
$VP$ (mm) & $D_1$ ($\mu$m) & $D_2$ ($\mu$m) & $D_3$ ($\mu$m) & $E_k$ ($\mu$m) \\
\midrule
50 & 0 & 0 & 0 & 0 \\
100 & 0 & 50 & 0 & 50 \\
150 & 0 & 50 & 50 & 50 \\
200 & 0 & 50 & 50 & 50 \\
250 & 0 & 0 & 0 & 0 \\
\bottomrule
\end{tabular}
\end{table}


Sustituyendo los valores de la Tabla~\ref{tab:pie-rey-diferencias}, la mayor
dispersión ($50\ \mu\text{m}$) aparece simultáneamente en los puntos de $100$ mm y
$200$ mm, de modo que el máximo de \eqref{eq:E} no es único. Conforme al criterio del
procedimiento, se selecciona el de mayor valor nominal:

\begin{equation}
  \therefore\quad E = 50\ \mu\text{m} \quad \text{en } 200 \text{ mm}.
  \label{eq:E-valor}
\end{equation}

Este resultado identifica además el punto crítico del instrumento para la prueba de
repetibilidad. Bajo el supuesto de una distribución rectangular de semiamplitud
$E/2$, la ecuación~\eqref{eq:sigma-rectangular} con $a = E/2$ proporciona

\begin{equation}
  \sigma(L_{i_E}) = \frac{2a}{2\sqrt{3}} = \frac{E}{2\sqrt{3}}
  = \frac{50}{2\sqrt{3}}\ \mu\text{m} \approx 14.43\ \mu\text{m}.
  \label{eq:sigma-LiE}
\end{equation}
